%%
%% SenSys '26 submission, acmart sigconf template.
%%
\documentclass[sigconf,anonymous,nonacm]{acmart}

\usepackage{makecell}
\usepackage{booktabs}
\usepackage{subcaption}
\usepackage{xcolor}

\newcommand{\edited}[1]{#1}
\newcommand{\scaffold}[1]{\textcolor{red}{#1}}

\AtBeginDocument{%
  \providecommand\BibTeX{{%
    Bib\TeX}}}

% \copyrightyear{2026}
% \acmYear{2026}
% \setcopyright{cc}
% \setcctype{by}
% \acmConference[SenSys '26]{ACM/IEEE International Conference on Embedded Artificial Intelligence and Sensing Systems}{May 11--14, 2026}{Saint Malo, France}
% \acmBooktitle{ACM/IEEE International Conference on Embedded Artificial Intelligence and Sensing Systems (SenSys '26), May 11--14, 2026, Saint Malo, France}
% \acmDOI{}
% \acmISBN{}

% System name is a placeholder until we settle on one.
\newcommand{\sys}{\textit{proposed system}}
\newcommand{\Sys}{\textit{Proposed system}}

% --- Paper vs. code terminology (codebase is NOT being renamed) -------------
%   Paper term           Code/filesystem term
%   -----------------    ---------------------------------------------------
%   proposed system  <-  "AVP" / "Active Video Perception"
%                        (e.g. system_design/06_avp_round1*.py, cache/avp*/,
%                        output files avp_*_preds.json, planner, etc.)
%   reasoner         <-  "observer" in code (the VLM amount-estimation step:
%                        run_observer(), OBSERVER_PROMPT, session_log.observer,
%                        --observer-only CLI flag, ...).
%   observer         <-  the low-level CV stack that builds the planner prompt:
%                        Hands23 (01), SigLIP (02), DINOv2 (03), OWLv2 scene
%                        (07d). No single identifier in code.
% Keep this block in sync with project_terminology_mapping.md in memory.
% ---------------------------------------------------------------------------

\pagestyle{plain}
\begin{document}

\title{Agentic Edge–Cloud Orchestration for Resource-Efficient \\ Daily Inventory Monitoring with Egocentric Video}

% \author{Kailai Cui}
% \affiliation{%
%   \institution{University of Michigan -- Ann Arbor}
%   \city{Ann Arbor}
%   \state{MI}
%   \country{USA}}
% \email{kailaic@umich.edu}

\begin{abstract}
We target automatic, per-item inventory tracking in home kitchens from
first-person video. Given a participant's current inventory (items with known
starting amounts) and a cloud-uploaded kitchen-activity video, the system
outputs, for each inventory item, how much was consumed and how much remains,
under an objective of \emph{minimizing VLM cost subject to an accuracy
 requirement}. 
 % The core insight is that inventory-relevant interactions occupy
% only a small fraction of cooking video (roughly 6\% in HD-EPIC), so lightweight
% CV models can pre-select the short observation windows over which an expensive
% VLM must reason. We are collecting a benchmark of in-the-wild and
% controlled-consumption recordings across four households with Meta Ray-Ban
% Smart Glasses, and we evaluate against per-session ground-truth consumption
% using a capacity-normalized percentage error (CNPE). This draft reports the
% problem statement, system design, data-collection targets and current
% progress, and the evaluation plan; introduction, related work, and results
% are deferred to later revisions.
\end{abstract}

% \begin{CCSXML}
% <ccs2012>
%    <concept>
%        <concept_id>10010147.10010178.10010224</concept_id>
%        <concept_desc>Computing methodologies~Computer vision</concept_desc>
%        <concept_significance>500</concept_significance>
%        </concept>
%    <concept>
%        <concept_id>10003120.10003138</concept_id>
%        <concept_desc>Human-centered computing~Ubiquitous and mobile computing</concept_desc>
%        <concept_significance>300</concept_significance>
%        </concept>
%  </ccs2012>
% \end{CCSXML}

% \ccsdesc[500]{Computing methodologies~Computer vision}
% \ccsdesc[300]{Human-centered computing~Ubiquitous and mobile computing}

% \keywords{egocentric video, food inventory, vision--language models, wearable sensing, active perception}

\maketitle


\section{Introduction}
% Motivation: Why people want the daily inventory monitoring.
% Motivation: How egocentric video captured by smartglass can help?
% Motivation: One example to give people an imagination about the usage scenario.
Wearable AI systems are emerging as a promising platform for augmenting human memory in everyday life. 
%
Imagine a user standing in a grocery store and asking their smartglasses, “Do I still have eggs at home?” 
%
Answering this question requires more than recognizing objects in the current scene: the assistant must recall prior observations from daily activities, retrieve the current household inventory, and reason about how that inventory has changed as items are purchased, stored, consumed, or discarded. 
%
Such capability could address common and consequential challenges in daily food management. 
%
% https://www.fda.gov/food/consumers/food-loss-and-waste
% https://www.cdc.gov/food-safety/data-research/facts-stats/index.html
In the United States, 30–40\% of the food supply is wasted, costing an average family of four approximately $1,500$ per year, while foodborne illness affects an estimated $48$ million people annually and causes $128,000$ hospitalizations and $3,000$ deaths \cite{}.

% Limitatioin of prior arts and bottleneck challenge.
% 1. without considering the cost of continous video understanding using VLM reasoning.
% 2. have to use advanced VLM with long-horizon contextual memory for understanding 
% 1. VLM reasoning cost for continuous video reocrding.
% 2. VLM context window for continuous video understanding.
% 3. Accuracy of the daily inventory monitoring.
% 4. Unstable egocentric view.
Recently, the cloud-based advanced large vision-language models (VLMs) and agentric AI has enabled powerful open-vocabulary visual understanding, making it possible to recognize diverse food items, infer object states, and answer natural-language queries about a scene \cite{tateno2025learning, zameni2025moscato,ukai2025status}.
%
However, enabling practical wearable memory assistance from egocentric video remains highly challenging.
%
A straightforward approach is to continuously stream smartglass video to a cloud VLM and ask it to summarize inventory-related observations, but such a design is difficult to deploy in daily life.
%
\emph{First}, continuous VLM reasoning incurs substantial inference cost.
%
\emph{Second}, inventory memory construction requires long-horizon reasoning to identify meaningful inventory updates.
%
Simply increasing the amount of video provided to a VLM is constrained by context-window limits.
%
\emph{Third}, egocentric video is inherently unstable, redundant, and incomplete, with frequent viewpoint changes, hand occlusions, and brief object appearances.
%
Consequently, continuous VLM reasoning requires a high-end, large-context VLM such as Gemini~3 Pro, yet still achieves limited inventory-monitoring accuracy.
%
At Gemini~3 Pro paid-tier rates~\cite{googleGeminiAPIPricing2026}, processing 1~fps egocentric video costs \$3.45 per hour, exceeding \$100 per household per month with only one hour of daily logging.

% Our solution:
% 1. selective. through the applied lens of automated kitchen food inventory tracking. 
% Not all the video require VLM reasoning depending on the targeted goal.
% 2. agentic edge–cloud orchestration
% 3. 
To resolve these challenges, our key insight is that not all egocentric video segments require expensive cloud VLM reasoning.
%
Only a small subset of moments involve potentially meaningful inventory changes, such as object acquisition, storage, removal, consumption, relocation, or state change.
%
This leads to the research question of this paper: can lightweight edge intelligence identify which egocentric video segments require VLM reasoning, thereby reducing overall token consumption while preserving or even improving daily inventory-monitoring accuracy?
%
This motivates a selective agentic architecture that uses lightweight edge intelligence to observe the video stream, identify candidate inventory-relevant events, maintain short-term context, and decide when cloud VLM reasoning is necessary.

We present \sys{}, an agentic edge--cloud orchestration framework for resource-efficient daily inventory monitoring using egocentric video.
%
Through the applied lens of kitchen food inventory monitoring, \sys{} integrates lightweight edge observers, agentic event selection, cloud VLM reasoning, and structured memory updates.
%
Specifically, smartglasses record egocentric video stream and upload it to a household edge endpoint, such as a smartphone, home desktop, or local server.
%
The edge endpoint first runs lightweight computer-vision observers over the full session, producing compact, timestamped logs of inventory-relevant interactions.
%
An LLM-powered agent then acts as a remote planner over these logs: it decides which short video windows warrant deeper reasoning, dispatches the selected segments to a cloud VLM reasoner, and updates a structured inventory ledger.



\section{Introduction}
\label{sec:intro}
The increasing ubiquity of wearable egocentric cameras enables continuous, passive memory augmentation for daily life logging and contextual assistance.
% In the egocentric vision community, this is formalized as Episodic Memory Retrieval [cite Ego4D], where systems must answer natural language queries regarding past events.
A critical component of these applications is the ability to track object state changes, as physical states reflect the outcomes of daily actions that humans frequently forget (e.g., determining whether a stove was left on, or quantifying how much of an ingredient was dispensed)[cite]. Because these queries require sophisticated causal and physical reasoning, recent works have increasingly relied on large vision--language models (VLMs) to interpret these dynamic transformations \cite{tateno2025learning, zameni2025moscato,ukai2025status}.

% Candidate names for this task class (defer naming to Related Work / Problem Statement):
%   1. "Selective episodic state retrieval" -- emphasizes selection + memory framing; closest to Ego4D NLQ vocabulary.
%   2. "Persistent inventory-style state tracking" -- emphasizes list-of-items setup; commits framing to inventory.
%   3. "Long-horizon state-change retrieval over a query set" -- most descriptive, least catchy.
Yet these VLM-based approaches have been validated predominantly in \emph{procedural settings} --- recipe execution, assembly, instruction following --- where every action advances the task and state changes saturate the timeline \cite{peddiCaptainCook4DDatasetUnderstanding2024, liBuildingEgocentricProcedural2025}. Daily memory augmentation, however, is fundamentally \emph{selective}: the user maintains a finite set of items of interest (a grocery ledger, a list of appliances, a checklist of medications) and asks about their state across long, unstructured recordings of ordinary life. In this regime, interactions with the queried items are sparse within hours of unrelated activity --- yet each interaction still demands the same heavy VLM reasoning that procedural settings rely on. This exposes a critical systems bottleneck: feeding hours of always-on video into VLMs incurs prohibitive API token costs \cite{openaiAPIPricing2026,googleGeminiAPIPricing2026} and far exceeds the context window of State-of-the-art VLMs \cite{yang2025qwen3, singh2025openaigpt5card, google2025gemini3promodelcard}, even though the task-relevant moments occupy a small fraction of the timeline. At Gemini~3 Pro paid-tier rates~\cite{googleGeminiAPIPricing2026}, end-to-end inference on 1~fps egocentric video costs \$3.45/hour --- over \$100/household/month at one hour of daily logging. How, then, can we achieve accurate, long-horizon state retrieval over a query set without exhausting the VLM token budgets that this reasoning consumes?

% Existing works on egocentric video understanding explores object state changes primarily through offline classification on trimmed video clips.
% moscato,xxx,assembly,cooking procedural, 
% However, another important yet unaddressed aspect of understanding object state change is the system efficiency on untrimmed long-form videos. [](what? inventory tracking? object retrieval?)
% Executing multimodal foundation models over always-on video incurs prohibitive computational costs (examples of gemini api cost), yet their deep reasoning is required for the brief moments when state changes actually occur. how can we optimize for computation cost efficiency with the constraint of accurate state change recall and understanding?

We answer this question through the applied lens of automated kitchen food inventory tracking. Food inventory is a near-universal daily memory-augmentation need: households replenish and consume food on a near-daily cycle, and the gap between what is on hand and what the user remembers drives wasted food, missed nutrition goals, and redundant grocery trips [cite]. It also introduces visual perception challenge  --- food morphs across fluids, solids, granulars, and discrete countables, each calling for a different combination of image retrieval, action recognition, container-aware detection, and volume- or count-based reasoning --- exactly the heavy VLM capabilities whose token cost is the bottleneck identified above. 
% The combination of high deployment frequency and a perceptually heterogeneous item set makes food inventory a stringent testbed for cost-efficient, long-horizon state retrieval. 
Concretely, given untrimmed egocentric video of cooking activities and an initial ledger, our goal is to quantify the consumed and remaining amount of each inventory item while minimizing the VLM token cost subject to \scaffold{accuracy requirement}. 

% --- DRAFT BULLETS for the rest of the introduction --------------------
% Logical order chosen here: (i) propose system at high level (one sentence,
% no methodology), (ii) state the two challenges that any such system must
% solve, (iii) describe how our design addresses them (CV pipeline +
% agentic plan-observe), (iv) dataset contribution, (v) contributions list.
% Bullets are intentionally raw -- to be turned into prose on the next pass.

We propose \sys{}, which separates inventory tracking into a \emph{locating} stage and a \emph{reasoning} stage. Smart glasses record kitchen video and upload it to a household edge endpoint --- a smartphone, a home desktop, or a local server. The endpoint first runs a lightweight CV pipeline locally over the full session, producing a compact, timestamped record of inventory interactions. It then runs an agent that calls a remote LLM planner over this record to decide which short segment windows are worth observing, dispatches those segments to a remote VLM observer, and uses the observer's per-item answers to update the inventory ledger. 

\paragraph{Challenges.}
\begin{itemize}
  \item \textbf{C1 -- Locating inventory-interaction windows in long video.} Kitchen sessions span tens of minutes to hours; only a few-percent slice involves actual contact with a ledger item. The system must surface those moments cheaply, before any heavy VLM reasoning is invoked.
  \item \textbf{C2 -- Selecting the few decision-relevant frames within those windows.} Even after C1, most located frames do not answer ``how much of item X remains?'' --- pre-interaction states, hand occlusions, wrong-item grabs, mid-pour blurs all coexist with the one or two frames that actually expose the post-interaction fill level. The system must read as few frames as possible while still landing on the answer-bearing ones.
\end{itemize}

\paragraph{Our design.}
\begin{itemize}
  \item Inspired by recent work on long egocentric video understanding [cite] and agentic video analytics [cite], we decompose the task into a \emph{lightweight CV pipeline} that builds a compact, item-indexed context of inventory-interaction events, followed by an \emph{LLM-driven agentic loop} that plans which short windows the VLM should observe, reads the VLM's per-item answers, and replans until each item's remaining amount is resolved or a token budget is exhausted.
  \item C1 is addressed by the CV pipeline: hand--object interaction detection, scene tagging, and item identification jointly produce a per-item temporal event list that locates candidate windows without invoking a VLM on the full video.
  \item C2 is addressed by the agentic plan--observe loop: a text-only LLM allocates a shared frame budget across items based on the CV evidence; a VLM observes only the planner-selected frames and returns per-item answers; the planner then replans on items whose state remained uncertain.
\end{itemize}

\paragraph{Dataset.}
\begin{itemize}
  \item Existing egocentric kitchen datasets annotate actions, not persistent item state, so we collect and annotate our own benchmark: in-the-wild and controlled-consumption recordings across multiple households on Meta Ray-Ban Smart Glasses, with per-session ground-truth starting and remaining amounts for every item the participant interacts with.
\end{itemize}

\paragraph{Contributions (placeholder list).}
\begin{itemize}
  \item A formulation of selective-query inventory tracking from untrimmed egocentric video as a cost--accuracy problem, with VLM token cost as the cost axis.
  \item \sys{}: a CV-located, agentically-planned VLM observation pipeline that addresses C1 and C2.
  \item A new benchmark of in-the-wild and controlled household kitchen video with per-item starting and remaining amounts.
  \item Cost--accuracy frontier evaluation against a whole-video VLM baseline on the new benchmark, with cross-cuts by item shape category and per-item lifecycle.
\end{itemize}
% --- END DRAFT BULLETS --------------------------------------------------
% Commented out: superseded by Section~\ref{sec:motivation} (Background and
% Motivation). Content folded into sec:bg-locate. Kept here for reference
% in case any wording is worth pulling forward.
%
% \subsection{Motivating observations}
%
% \begin{itemize}
%   \item \textbf{Videos are long; the relevant window is short.} Updating an
%   item's inventory only requires observing the brief interaction with that
%   item, not the rest of the session.
%   \item \textbf{Amount reasoning is hard and needs a powerful VLM.} Deciding
%   how much of a jar of sauce was used, or how many tomatoes were cut, is a
%   fine-grained visual-reasoning task that weaker models do poorly on.
%   \item \textbf{Powerful VLM APIs are expensive.} An hour of video sampled at
%   1~fps costs roughly \$2.5 in Gemini~3~Pro input tokens alone, and current
%   VLM context windows typically cannot fit more than about 10~minutes of
%   video per call.
%   \item \textbf{Most of the video is irrelevant.} In HD-EPIC, interactions
%   with inventory items account for only $\sim$6\% of cooking video time; the
%   rest is meal prep, cooking, washing dishes, and similar activity
%   \cite{perrettHDEPICHighlyDetailedEgocentric2025}.
% \end{itemize}
%
% \subsection{Design insight}
%
% Use lightweight CV models to \emph{locate} segments that interact with
% inventory items, so that an expensive VLM only reasons over those short
% windows. This both reduces token cost and makes it feasible to reason across
% long-horizon recordings that would otherwise exceed the VLM's context limit.

\section{Background and Motivation}
\label{sec:motivation}

\subsection{Related work}
\label{sec:bg-related}

\noindent\textbf{Long-form egocentric video understanding.}
Long-horizon egocentric video understanding has been driven by
progressively larger and more densely annotated datasets.
Ego4D~\scaffold{[cite]},
EPIC-Kitchens~\cite{damenRescalingEgocentricVision2022}, and
HD-EPIC~\cite{perrettHDEPICHighlyDetailedEgocentric2025} together
supply thousands of hours of first-person activity labeled at
second-scale resolution with verb--noun action segments, narrations,
and hand--object interactions. Built on these substrates, recent
agentic VLM systems answer free-form queries over long video by
iteratively retrieving the few answer-bearing
frames~\scaffold{[cite VideoAgent, VideoTree, LLoVi, SeViLA]},
decomposing what would be a single long-context VLM call into a
planner--observer loop. We adopt the same locate-then-read
decomposition with two adaptations to our setting: the planner is
anchored to a fixed query set (the inventory ledger) rather than a
free-form natural-language question, and the observer is asked to
commit to a numerical amount per item rather than retrieve a span.

\noindent\textbf{Object state-change understanding.}
A separate strand formalises object transformations as classification,
recognition, or graph-prediction
targets~\cite{tateno2025learning, zameni2025moscato, ukai2025status}.
These approaches share our goal of recovering per-item state from
interaction, but evaluate predominantly on trimmed clips where the
relevant interaction is already isolated; they do not confront the
temporal selection problem that arises when the input is hours of
unstructured cooking video. \sys{} inherits the per-item state target
but operates on the long-video locating regime above.

\noindent\textbf{Existing kitchen datasets are action-centric.}
Quantifying either the sparsity or the displacement properties below
requires ground-truth per-item starting and remaining amounts on
untrimmed long-horizon recordings --- supervision that no existing
benchmark provides. The large-scale egocentric kitchen video datasets
cited above annotate \emph{actions} (verb--noun pairs, action
segments, narrations) but not the persistent \emph{state} of
inventory items across sessions. We address this gap with a new
dataset (\S\ref{sec:dataset}).

\subsection{Motivation study}
\label{sec:bg-motivation}

\noindent\textbf{The cost of running VLMs on whole videos.}
Naively reading whole sessions into a state-of-the-art VLM is the
obvious baseline and a useful cost reference. On a 10-session pilot
from our dataset ($\approx$2.2 hours of 1~fps egocentric cooking
video, 7{,}977 frames sent), end-to-end inference with
Gemini~3.1~Pro at full visual resolution costs \$7.65 at current
paid-tier rates~\cite{googleGeminiAPIPricing2026} --- \$3.45 per
hour of recording, exceeding \$100/household/month at one hour of
daily logging. The fundamental cost driver is the linear scaling of
input tokens in video
length~\cite{openaiAPIPricing2026,googleGeminiAPIPricing2026};
against that scaling, the relevant evidence we want the VLM to reason
about --- short interactions with inventory items --- is a small
fraction of the input.

\noindent\textbf{Inventory-relevant interactions are sparse.}
Cooking sessions are long, routinely tens of minutes to hours, while
the moments that update the ledger --- retrieval, dispensing,
put-back --- occupy only a small fraction of the timeline.
\scaffold{On HD-EPIC, manipulations of inventory-relevant items
account for roughly $X$\% of session duration
(Fig.~\ref{fig:sparsity}); EPIC-Kitchens shows comparable
density~\cite{damenRescalingEgocentricVision2022,perrettHDEPICHighlyDetailedEgocentric2025}.}
% Note: we frame this as 'selective' (a query set of inventory items
% in otherwise unstructured cooking) rather than 'procedural' (every
% action advances a single recipe, e.g. CaptainCook4D, Assembly101).
This sparsity is what makes locate-then-read viable: lightweight CV
signals (HOI contact, item identification, scene context) can
pre-select the seconds in which an item is actually being touched, so
an expensive VLM is only invoked on planner-allocated frames rather
than the whole session.

\noindent\textbf{The inventory-relevant moment is often displaced from the obvious interaction.}
Sparsity tells us \emph{where to look first} but not \emph{which
frames contain the answer}. An HOI detection identifies the seconds
when the user is in contact with an item, yet the amount-relevant
evidence --- the moment a bottle is decanted, an egg is broken, a bag
is opened --- is often offset from the contact peak: it can occur
seconds later in the same window, on a return trip after an
interleaved sub-task, or behind a hand. The inventory-changing event
may even be split across the session (pickup at minute~3, decant at
minute~8), so a single tight window around the contact peak misses
the answer entirely.\scaffold{ [cite agentic / iterative
frame-retrieval works that converge on the same observation in
long-video QA: VideoAgent, VideoTree, LLoVi, SeViLA, Goldfish].} The
challenge, then, is to budget the VLM's frame reads across both the
contact peak and its temporal context, and to revise that budget
when the first read finds no answer --- exactly what motivates the
two-stage planner in \S\ref{sec:system}.

\section{Dataset Overview and Collection Progress}
\label{sec:dataset}

The dataset is collected from three participants cooking in their own
kitchens. Each participant records first-person video using either
\textbf{Meta Ray-Ban Smart Glasses} (egocentric, fisheye) or a head-mounted
\textbf{Raspberry Pi Camera Module~v3}, depending on hardware availability.
For each session we also record, for every item the participant interacts
with, the item's starting amount and ending amount (by weight for bulk
items, by count for discrete items such as eggs), which serve as the
ground-truth consumption label.

\subsection{Collection protocol}

We collect two complementary flavours of data with the same recording
hardware and the same start/end measurement procedure.

\paragraph{In-the-wild cooking sessions.} Participants wear the recording
device during their ordinary meal-prep sessions at home, with no
instructions on what or how to cook. These sessions cover the realistic
target scenario --- free-form cooking with multiple items, varied
background clutter, and whatever item subset happens to be on the shelf
that week --- and form the primary evaluation cohort.

\paragraph{Controlled consumption sessions.} In a controlled session the
participant follows a scripted protocol that consumes a \emph{single}
pre-selected item at a time: the item is removed from storage, a known
amount is poured/scooped/cut out, and the item is returned. Several items
are cycled through in a single recording. We include controlled sessions
for two reasons:
\begin{enumerate}
  \item \textbf{Food variety.} In-the-wild cooking sessions naturally
  concentrate on a small set of staples that a given participant happens to
  cook with; controlled sessions let us deliberately exercise items that
  are under-represented in free-form cooking (e.g.~less-used condiments,
  opaque-packaged items, items with unusual shape categories), broadening
  the visual and packaging distribution the evaluation covers.
  \item \textbf{Statistical power.} Controlled sessions produce many
  cleanly-bounded consumption events per unit recording time, which
  increases the event count available for paired accuracy comparisons and
  allows us to detect \sys{}-vs.-baseline improvements in VLM
  amount-estimation accuracy at statistical significance that the wild
  cohort alone, with its smaller event count, would not reliably support.
\end{enumerate}

\subsection{Current collection status}
\label{sec:current}

Table~\ref{tab:progress} reports the dataset as it currently stands
against the targets for the full benchmark (4 households, 20~h in-the-wild
+ 2~h controlled, 400 + 300 consumption events, $\sim$70 distinct items
across 4 shape categories). Four participants have begun recording; three
(\texttt{kailai}, \texttt{giraffe}, \texttt{chaoyu}) have weighed-in
consumption events logged to a ledger. A fourth participant's ledger is
populated with purchases only (annotation in progress); a fifth has
started recording but has not yet begun annotation.

\begin{table}[h]
\centering
\caption{Current collection progress vs. target.}
\label{tab:progress}
\begin{tabular}{lrr}
\toprule
                                 & Current & Target \\
\midrule
Participants (recording)         & 5       & 4      \\
Participants (annotated ledger)  & 3       & 4      \\
Total video                      & 28.0~h  & 22~h   \\
\quad in-the-wild                & 27.5~h  & 20~h   \\
\quad controlled                 & 0.4~h   & 2~h    \\
Consumption events (weighed)     & 409     & 700    \\
\quad in-the-wild                & 377     & 400    \\
\quad controlled                 & 32      & 300    \\
Distinct tracked item instances  & 136     & $\sim$70 \\
\bottomrule
\end{tabular}
\end{table}

Across the annotated ledgers, the 136 tracked item instances fall into
four shape categories (Table~\ref{tab:items} below breaks this down with
examples): \textsc{solid} 39, \textsc{discrete} 37, \textsc{granular} 33,
\textsc{fluid} 20, with 7 instances not yet categorised. In-the-wild
hours already exceed the 20~h target, and the wild event count (377) is
close to its 400-event target. The controlled cohort is the principal
bottleneck --- only $\sim$0.4~h recorded and 32 events logged against
targets of 2~h and 300 events --- and is our near-term collection focus
(see rationale in Section~\ref{sec:dataset}).

\subsection{Item variety}
\label{sec:items}

Each item in the dataset is tagged with a \emph{shape category} --- one of
\{\textsc{fluid}, \textsc{granular}, \textsc{solid}, \textsc{discrete}\}
--- along with a \emph{transparency} flag indicating whether the package
is see-through. These attributes are informative because they govern how
an amount is estimated: fluid and granular items in transparent packaging
permit fill-level reading, discrete items (eggs, yogurt cups) permit
counting, and opaque containers must be inferred from pour/scoop duration
or weighing context. Table~\ref{tab:items} summarises the current item
inventory.

\begin{table}[h]
\centering
\caption{Item variety across all annotated ledgers, grouped by shape
category. Counts are item \emph{instances} (one row per purchase).
Transparency labelling is ongoing.}
\label{tab:items}
\begin{tabular}{lrl}
\toprule
Category    & Items & Examples \\
\midrule
Fluid       & 20 & olive oil, milk, soy sauce \\
Granular    & 33 & rice, pasta, cereal \\
Solid       & 39 & garlic, steak, tofu \\
Discrete    & 37 & eggs, yogurt cups, bagels \\
Unlabelled  & 7  & (categorisation pending) \\
\midrule
Total       & 136 & \\
\bottomrule
\end{tabular}
\end{table}

\subsection{Longitudinal coverage}
\label{sec:longitudinal}

The same item often persists across many sessions before it is exhausted,
which lets us evaluate \emph{how a system tracks one package over its
full lifecycle} rather than only how it does on isolated sessions.
Across the currently-annotated ledgers, item lifecycles span the
distribution shown in Table~\ref{tab:longitudinal}: 32 item-instances
appear in at least 5 distinct sessions, 11 in at least 10, and the
longest-running items (eggs, sliced cheese, cereal, milk gallons) are
observed across 13--18 sessions each. This long-tail subset is the
substrate for the per-item lifecycle evaluation described in
Section~\ref{sec:lifecycle}.

\begin{table}[h]
\centering
\caption{Item lifecycle distribution across annotated ledgers.
``Sessions'' is the number of distinct sessions in which a given item
instance was logged as consumed.}
\label{tab:longitudinal}
\begin{tabular}{rrr}
\toprule
Sessions per item & Item instances & Cumulative ($\geq$) \\
\midrule
$\geq$1   & 92 & 92 \\
$\geq$3   & 55 & 55 \\
$\geq$5   & 32 & 32 \\
$\geq$10  & 11 & 11 \\
$\geq$15  & 3  & 3  \\
\bottomrule
\end{tabular}
\end{table}

The bulk of this longitudinal coverage currently comes from a single
participant (\texttt{kailai}, 70 items). Bringing the other participants
to comparable lifecycle depth, rather than to a wider item catalogue, is
the principal collection focus alongside the controlled-cohort gap noted
in Section~\ref{sec:current}.

\section{System Design}
\label{sec:system}

\sys{} takes an untrimmed egocentric kitchen video and a per-participant
\emph{ledger} (the list of purchased food items with package amounts and
reference images), and produces a per-session update of \emph{how much of
each item was consumed}. The pipeline is organized as a sequence of
numbered stages (scripts in \texttt{system\_design/}) with well-defined
inputs and outputs so that components can be swapped or ablated.

\subsection{Inputs and assumptions}
\label{sec:system-inputs}

We assume that each item is registered \emph{at the time of purchase}
with two lightweight artifacts:
\begin{itemize}
  \item A \textbf{receipt entry}, supplying the item's identity,
  quantity, and unit (grams or count). This is what initialises the
  per-participant ledger; subsequent sessions update an item's remaining
  amount until it is exhausted.
  \item One or more \textbf{reference images} of the item's package,
  captured by the participant or pulled from the retailer / web. These
  are what the item-identification stage (DINOv3, below) matches the
  in-contact object crops against.
\end{itemize}
Both artifacts are produced once, when the item enters the household,
and reused across every session that touches that item. Within a
session, the system assumes nothing else: the video is processed
independently against the current ledger snapshot.

\subsection{Pipeline stages}

\begin{enumerate}
  \item \textbf{Hand--object interaction (HOI grasp detection).} The
  Hands23 detector runs at 1~fps across the full video and flags frames
  in which a hand is in \emph{object-contact} with something it is
  holding. Only these grasp frames are carried forward; frames with no
  active grasp are dropped, so downstream stages operate on a sparse
  set of moments when the participant is actually manipulating an item.
  \item \textbf{Scene tagging (OWLv2).} An OWLv2 open-vocabulary
  detector labels each grasp frame with a coarse kitchen location:
  \texttt{storage}, \texttt{sink}, or \texttt{stove}. These tags give
  the planner the lifecycle cue needed to distinguish retrieval from
  use from put-back when allocating observation budget.
  \item \textbf{Item identification (DINOv3).} The in-contact object
  crop from each grasp frame is matched against the participant's
  current inventory via DINOv3 image-to-image similarity against each
  item's reference images (Section~\ref{sec:system-inputs}). The
  matching score decides \emph{which ledger item} is being handled in
  each grasp frame. A SigLIP2 zero-shot text-to-image score is computed
  in parallel but is disabled in the canonical configuration, so the
  planner's evidence rendering relies on HOI contact and DINOv3 alone.
  \item \textbf{Amount estimation.} \sys{} predicts \emph{how much
  remained} (or how much was used) for each touched item at the end of
  the session. This stage is decomposed so that the expensive
  vision--language reasoning is invoked at most twice per session, each
  time on a planner-selected frame set:
  \begin{enumerate}
    \item[4a.] \textbf{Frame-budget planner (text-only LLM).} The
    grasp, identification, and scene evidence from stages~1--3 is
    collapsed via heuristic morphological compression (close-then-open
    on per-frame DINOv3 scores) into per-item, HOI-gated temporal
    segments --- a chronological list of on-contact bursts, each tagged
    with peak DINOv3 score and a coarse scene distribution. A text-only
    LLM reads the session inventory together with this evidence and
    emits, per inventory item, an \texttt{observe} /
    \texttt{no\_observation} decision plus a single shared list of
    sweep timestamps, partitioned into (i)~\emph{journey samples}
    spread across the item's bursts and (ii)~\emph{dense windows}
    concentrated on the most informative bursts (with a configurable
    rule biasing toward late, dispense-phase bursts over early
    retrieval). Each timestamp is tagged with the candidate items the
    observer should look for there, so multiple items can be searched
    for in a single observer call.
    \item[4b.] \textbf{Sweep VLM observer.} A \emph{single}
    multi-image VLM call sees the merged sweep frames and returns, for
    every candidate item the planner flagged on those frames,
    \texttt{\{item\_confirmed, evidence\_frames, amount\_remaining\}}.
    The observer is multi-item by design (one call per round, not one
    call per item), so the per-frame token cost is amortised across
    every item the planner asked about on that frame. Amount estimation
    --- judging how much of a jar of sauce was poured, or how many
    tomatoes have been removed from a bag --- demands fine-grained
    visual reasoning that weaker models still handle poorly
    \scaffold{[cite MMVP, BLINK, V*]}; we therefore pay for a strong
    VLM at this stage (Gemini 3.1 Pro Preview in our canonical
    configuration~\cite{googleGeminiAPIPricing2026}) and amortise its
    per-call cost across all items flagged for the same sweep frames.
  \end{enumerate}
\end{enumerate}

Stages~1--3 correspond to the lightweight ``locate relevant windows''
path described in Section~\ref{sec:motivation}; stage~4a is where a
cheap language model reasons over the full session at the symbol
level, and stage~4b is where the powerful VLM performs the expensive
perceptual reasoning, but only on the planner-allocated frames.

\paragraph{Plan--observe--replan (two rounds).} Stage~4 is structured
as a fixed two-round plan--observe--replan loop. After round~1, the
planner is re-invoked with (i)~the actual timestamp ranges of frames
the round-1 observer saw (not the planner's declared windows) and
(ii)~the per-item answers and confidences the observer returned. The
round-2 planner is allowed to revisit any moment in the session ---
including frames already inspected --- and dispatches a second sweep
observer call targeting items whose remaining amount, starting amount,
or per-event derivative was unresolved or under-confident in round~1.
The loop terminates after round~2, or earlier if the planner declares
all items resolved.

\scaffold{\subsection{Reasoner extensions}}
\label{sec:extensions}

\scaffold{\paragraph{Cross-session visual memory.} Amount estimation
within a single session is fundamentally under-determined when the stock
container is only partially visible, or when it never appears in its final
state. This extension threads a small \emph{visual memory} across
sessions: for each item the reasoner confirms, we persist a single
canonical end-of-session frame (the moment the stock container is most
clearly visible) together with its estimated remaining amount. On
subsequent sessions the reasoner is prompted with this item's prior
snapshots, in chronological order, alongside the current-session frames.
The prior \emph{numeric} estimates are withheld to avoid anchoring; only
the visual evolution of the package is exposed. The intended effect is
that the reasoner can compare the current fill level against the previous
fill level directly, rather than re-deriving the remaining amount from
scratch each session.}

\subsection{Baseline}

A \textbf{whole-video} baseline gives the VLM the entire untrimmed video
and asks it to predict remaining amounts end-to-end, with no temporal
structure supplied. This is the comparator against which the \sys{}'s
cost--accuracy gain is measured.

\section{Benchmark and Evaluation Plan}
\label{sec:evaluation}

\subsection{Benchmark structure}

Videos are organized into \emph{sessions} (meal-prep episodes). For each
session, the ground-truth label is, for every item touched, the
\emph{starting amount} and the \emph{end amount}. The system under test takes
a session video as input and produces a list of
\texttt{[item, remaining\_amount]} pairs as output.

\subsection{Evaluation cohorts}
\label{sec:eval-cohorts}

We evaluate the two collection flavours for different purposes.

\paragraph{In-the-wild cooking videos.} These are our primary cohort and
the one that reflects the deployed use case. On this cohort we report
\emph{both} (i) amount-estimation accuracy (CNPE, Section~\ref{sec:cnpe})
and (ii) VLM token consumption per session
(Section~\ref{sec:tokens}). Target scale at submission:
$\sim$20~h of video and $\sim$400 consumption events across 3
participants (numbers to be finalised).

\paragraph{Controlled consumption videos.} We use the controlled cohort to
evaluate VLM amount-estimation accuracy over a broader, deliberately
diversified set of food items (see Section~\ref{sec:dataset}). The item
mix here spans shape categories and packaging transparencies that are
under-represented in any one participant's natural cooking, which lets us
measure how \sys{}'s accuracy gain over baselines generalises across item
types. Target scale: $\sim$2~h of video and $\sim$300 consumption events;
we report CNPE on this cohort with paired per-event comparisons so that
\sys{}-vs.-baseline differences can be tested at statistical significance.

\subsection{Primary metric: CNPE (accuracy)}
\label{sec:cnpe}

We measure accuracy with \textbf{Capacity-Normalized Percentage Error (CNPE)}:
\[
\mathrm{CNPE} \;=\; \frac{\lvert \mathrm{prediction} - \mathrm{GT} \rvert}{W_0} \times 100,
\]
where $W_0$ is the item's original package size (capacity). Normalizing by
$W_0$ makes errors comparable across items of very different sizes: a 50~g
error on a 500~g jar (10\% CNPE) is treated the same as a 50~g error on a
50~g packet (100\% CNPE). CNPE can be computed equivalently on the remaining
amount or the used amount.

\subsection{Primary metric: VLM token consumption (cost)}
\label{sec:tokens}

We report total VLM token cost per session (input~+~output, with
thinking tokens charged at output rate). Cost is treated on equal
footing with CNPE: the central claim of the paper is a
\emph{cost--accuracy} trade-off, and a system that wins on one axis
while losing on the other is not preferred unless the trade is
explicitly characterized along the cost--accuracy frontier.

\subsection{Cross-cut analyses}
\label{sec:crosscut}

The two headline metrics (CNPE, token cost) are reported as session-level
aggregates, but the questions that drive system design are sharper than a
single mean. We additionally report three cross-cut breakdowns of the
same underlying per-event predictions.

\paragraph{Per-class CNPE.} CNPE is broken down by item shape category
(\textsc{fluid}, \textsc{granular}, \textsc{solid}, \textsc{discrete})
and crossed with packaging transparency. The hypothesis is that the
\sys{}'s gain over the baseline concentrates on classes where amount
reasoning is locally hard but \emph{visible} given the right frame ---
fill levels in transparent fluids, residue in granulars, missing items
in discrete counts --- whereas opaque-fluid and irregular-solid
estimation remain hard regardless of how the observation windows are
chosen. Reporting per-class CNPE makes the source of any aggregate
gain auditable and exposes packaging types where future work is most
needed. The per-class event counts in Table~\ref{tab:items} (20--39
instances per category, with multiple consumption events per instance)
are sufficient to support per-class means with bootstrap intervals; the
controlled-cohort expansion (Section~\ref{sec:current}) is targeted at
the categories currently thinnest in event count.

\paragraph{Item-lifecycle CNPE.}
\label{sec:lifecycle}
Inventory tracking is intrinsically a longitudinal task: the same
package is observed across many sessions until exhausted. For the
long-tail item subset characterised in Section~\ref{sec:longitudinal}
(32 item-instances with $\geq$5 sessions, 11 with $\geq$10), we report a
per-item CNPE \emph{trajectory} indexed by session number, alongside a
single per-item lifecycle CNPE (mean across that item's sessions). Two
trajectory shapes matter and we will distinguish them: (i)
session-by-session error that stays roughly flat, indicating the system
re-anchors against the current frame and earlier session errors do not
compound; and (ii) error that grows monotonically as the package
empties, indicating that low-fill states are a systematic failure mode
and motivating the cross-session visual memory extension
(Section~\ref{sec:extensions}). For items observed all the way to a
``finished'' event, we additionally report end-of-life cumulative drift
(predicted total used minus actual, normalised by capacity).

\subsection{Experiment matrix}

Every cell in the matrix below is reported under both cross-cut
breakdowns from Section~\ref{sec:crosscut} (per-class and
item-lifecycle), so that the headline CNPE/token comparisons can be
audited along the dimensions most likely to drive system design.

\begin{itemize}
  \item \textbf{Controlled vs.\ natural sessions.} We expect
  amount-estimation error to be lower on controlled sessions (single item,
  scripted consumption) and will report the two cohorts separately.
  \item \textbf{Baseline vs.\ \sys.} Whole-video baseline and the
  \sys{} (planner~+~sweep observer with iterative replan per
  Section~\ref{sec:system}), using the same VLM backends for a fair
  comparison.
  \item \textbf{Ablations within the \sys.} (i) Number of plan--observe
  rounds (1 vs.\ 2 vs.\ unbounded under a session-level frame budget);
  (ii) frame-budget allocation policy (journey-only, dense-only,
  mixed; late-burst-biased vs.\ uniform); (iii) planner evidence
  rendering (per-item temporal segments vs.\ chronological timeline
  vs.\ activity blocks); (iv) backbone VLM for the observer (GPT-5,
  Gemini~3 Pro / 2.5 Pro / Flash, Gemma) under matched frame-budget
  caps so the comparison isolates VLM choice from frame count.
  \scaffold{\item \textbf{Extension (Section~\ref{sec:extensions}).} We
  will additionally report \emph{cross-session visual memory}: the
  \sys{} with prior end-of-session snapshots threaded into the reasoner
  prompt, versus the same \sys{} without --- evaluated specifically on
  the long-tail items of Table~\ref{tab:longitudinal} where the
  trajectory metric (Section~\ref{sec:lifecycle}) makes any
  cross-session benefit visible. The extension is evaluated on the same
  CNPE/token axes so its cost/accuracy trade-off against the base
  \sys{} is directly comparable.}
  \item \textbf{Scalability.} As more participants and sessions come online,
  we will re-run the full matrix and report metric stability as the
  experiment set grows.
\end{itemize}

\subsection{Locked evaluation definitions}
\label{sec:lockdown}

To make the small-scale pilots and the final-cohort runs directly
comparable, three definitions are pinned for the remainder of the
project. Any later change to these is treated as a methodological
revision and the affected runs are re-executed.

\paragraph{Frozen cohort lists.} The wild and controlled cohorts are
specified as explicit session-ID lists per participant
(Appendix~\ref{app:cohort-ids}, to be added). New sessions recorded
after the freeze date are excluded from the headline numbers and may
appear only as out-of-distribution diagnostics. The current pilot
cohort is the 10-session canonical \texttt{kailai} natural-cooking
subset; the final cohort substitutes the larger frozen list under the
same protocol.

\paragraph{Frozen pricing convention.} Billable cost per session is
reported in USD using Gemini~3~Pro paid-tier pricing
(\$1.25/M~input,~\$10/M~output, thinking tokens charged at output
rate), with images converted at 1~frame~$=$~260 input tokens. The same
pricing convention is applied to the whole-video baseline and the
\sys{} so that cost comparisons are not confounded by per-comparator
pricing choices. Planner LLM cost is included only in the \sys{} and
extension rows.

\paragraph{Frozen metric definitions.} Three metrics constitute the
headline:
\begin{itemize}
  \item \textbf{CNPE per event} (Section~\ref{sec:cnpe}), reported as
  the mean over GT consumption events.
  \item \textbf{Recall} = the fraction of GT consumption events for
  which the system produced a non-null amount estimate. Items the
  planner explicitly skipped, items the observer marked as
  not-confirmed, and items never proposed in the first place all count
  as missed.
  \item \textbf{Fused CNPE${+}$recall} = the per-event mean of CNPE
  with missed events charged at $\mathrm{CNPE}=100$. This single
  scalar is the y-axis of the headline cost--accuracy figure. The
  failure floor of $100$ corresponds to a full-package error and
  prevents systems from improving the headline by skipping uncertain
  items.
\end{itemize}

\section{Results}
\label{sec:results}

\scaffold{This section is a skeleton. Each figure and table below is
defined here so that experiments executed between now and the
freeze date populate specific cells rather than producing orphaned
runs. Figures and tables will be filled as experiments complete.}

\subsection{Setup}
\label{sec:results-setup}

\scaffold{Final cohort sizes (sessions, hours, consumption events) per
participant, broken out by wild vs.\ controlled. One row per
comparator (whole-video baseline, \sys{} R1, \sys{} R1$+$R2,
\sys{}$+$memory) listing the locked run configuration: planner LLM,
observer VLM, frame-budget cap, evidence rendering mode,
burst-coverage rule, thinking depth. Pricing convention restated from
Section~\ref{sec:lockdown}.}

\subsection{Headline: cost--accuracy frontier}
\label{sec:results-headline}

\scaffold{\textbf{Figure (central):} Pareto plot. x-axis = USD/session
(log scale); y-axis = fused CNPE$+$recall (lower is better). One
marker per (comparator $\times$ frame-budget) cell. Curves: whole-video
baseline (single point), \sys{} frontier (sweep over frame budget and
round count). Iso-cost guides at \$0.10, \$0.50, \$1.00. Extension
points appended in red when implemented.}

\scaffold{\textbf{Table~\ref{tab:headline} (placeholder):} aggregate
CNPE, recall, and \$/session for each comparator on the wild and
controlled cohorts.}

\subsection{Cross-cut analyses}
\label{sec:results-crosscut}

\scaffold{Two panels populate the breakdowns specified in
Section~\ref{sec:crosscut}:}
\begin{itemize}
  \item \scaffold{\textbf{Per-class CNPE} bars: shape category
  (\textsc{fluid}/\textsc{granular}/\textsc{solid}/\textsc{discrete})
  $\times$ comparator, crossed with packaging transparency. Tests the
  ``gain concentrates on visible-but-locally-hard classes'' hypothesis
  from Section~\ref{sec:crosscut}.}
  \item \scaffold{\textbf{Lifecycle trajectory grid:} small multiples,
  one panel per long-tail item from Table~\ref{tab:longitudinal}, $x$
  $=$ session index, $y$ $=$ CNPE, two lines per panel (whole-video
  baseline / \sys{}). End-of-life cumulative drift reported in an
  adjacent table for items observed to ``finished''.}
\end{itemize}

\subsection{Ablations}
\label{sec:results-ablations}

\scaffold{Each ablation produces a small Pareto subplot sharing axes
with the headline figure. (i)~Round count (R1 vs.\ R1$+$R2 vs.\
unbounded under a session-level frame budget). (ii)~Frame budget
allocation policy (journey-only / dense-only / mixed; late-burst vs.\
uniform). (iii)~Planner evidence rendering (\texttt{blocks} vs.\
\texttt{chrono} vs.\ \texttt{segments}). (iv)~Observer-backbone swap
(GPT-5, Gemini~3~Pro, Gemini~2.5~Pro/Flash, Gemma) at matched
frame-budget caps.}

\subsection{Failure analysis}
\label{sec:results-failures}

\scaffold{2--3 hand-picked sessions with figure-strip illustrations:
session timeline, frames the observer actually saw, predicted vs.\
true amount, one-sentence diagnosis. Each case is mapped to either the
cross-session memory extension or flagged as an open problem.}

\scaffold{\subsection{Extension}}
\label{sec:results-extensions}

\scaffold{Cross-session visual memory results: matched-frame-budget
comparison against base \sys{}, evaluated on the
Table~\ref{tab:longitudinal} long-tail items so the lifecycle
trajectories from Section~\ref{sec:results-crosscut} can show the
benefit.}

\begin{table}[h]
\centering
\caption{\scaffold{(Placeholder.)} Headline CNPE, recall, and
\$/session per comparator on the frozen cohorts.}
\label{tab:headline}
\begin{tabular}{lrrrrrr}
\toprule
              & \multicolumn{3}{c}{Wild} & \multicolumn{3}{c}{Controlled} \\
\cmidrule(lr){2-4}\cmidrule(lr){5-7}
Comparator    & CNPE & Recall & \$/sess & CNPE & Recall & \$/sess \\
\midrule
Whole-video baseline & -- & -- & -- & -- & -- & -- \\
\sys{} (R1)          & -- & -- & -- & -- & -- & -- \\
\sys{} (R1$+$R2)     & -- & -- & -- & -- & -- & -- \\
\scaffold{\sys{}$+$mem} & -- & -- & -- & -- & -- & -- \\
\bottomrule
\end{tabular}
\end{table}

\section{Related Work}
long video understanding
ego: anchored by HOI and scene \cite{golettoAMEGOActiveMemory2024}, \cite{huang2025building}
agentic: avp, ava, 

object state change

Datasets/benchmark on Object States: follow MASCATO

\cite{wang2025active}

\bibliographystyle{ACM-Reference-Format}
\bibliography{references}

\end{document}
